% Part: computability
% Chapter: recursive-functions
% Section: primes

\documentclass[../../../include/open-logic-section]{subfiles}

\begin{document}

\olfileid{cmp}{rec}{pri}
\olsection{الأعداد الأولية}

بالتكميم المحدود والتصغير المحدود يتيسر لنا إثبات العودية البدائية لكثير من الدوال
والعلائق المألوفة. فلننظر، مثلًا، في العلاقة «$\OLId{latin-ordinary}{x}$ يقسم
$\OLId{latin-ordinary}{y}$»، التي تكتب $\OLId{latin-ordinary}{x}\mid \OLId{latin-ordinary}{y}$. ومعنى $\OLId{latin-ordinary}{x}\mid \OLId{latin-ordinary}{y}$ أن تقبل قسمة
$\OLId{latin-ordinary}{y}$ على~$\OLId{latin-ordinary}{x}$ دون باق، أي أن يكون $\OLId{latin-ordinary}{y}$ مضاعفًا صحيحًا لـ~$\OLId{latin-ordinary}{x}$. (وإذا لم
تصدق، أي إذا كان باقي قسمة $\OLId{latin-ordinary}{y}$ على $\OLId{latin-ordinary}{x}$ أكبر من~$\OLMathNumeral{0}$، كتبنا
$\OLId{latin-ordinary}{x}\nmid \OLId{latin-ordinary}{y}$.) وبعبارة أخرى، $\OLId{latin-ordinary}{x}\mid \OLId{latin-ordinary}{y}$ إذا وفقط إذا وجد~$\OLId{latin-ordinary}{z}$ بحيث
$\OLId{latin-ordinary}{x}\cdot \OLId{latin-ordinary}{z}=\OLId{latin-ordinary}{y}$. ومتى وجد شاهد $\OLId{latin-ordinary}{z}$ لهذا، أمكن اختيار شاهد يكون
$\leq \OLId{latin-ordinary}{y}$؛ حتى في حالة كون العددين صفرًا نختار الصفر شاهدًا. ومن ثم يصدق $\OLId{latin-ordinary}{x}\mid \OLId{latin-ordinary}{y}$ إذا وفقط إذا وجد $\OLId{latin-ordinary}{z}\le \OLId{latin-ordinary}{y}$ بحيث
$\OLId{latin-ordinary}{x}\cdot \OLId{latin-ordinary}{z}=\OLId{latin-ordinary}{y}$. فتعرّف العلاقة $\OLId{latin-ordinary}{x}\mid \OLId{latin-ordinary}{y}$ بالتكميم الوجودي
المحدود من المساواة والضرب على الوجه الآتي:
\[
\OLId{latin-ordinary}{x} \mid \OLId{latin-ordinary}{y} \defiff \bexists{\OLId{latin-ordinary}{z} \leq \OLId{latin-ordinary}{y}}{(\OLId{latin-ordinary}{x} \cdot \OLId{latin-ordinary}{z}) = \OLId{latin-ordinary}{y}}.
\]
فتثبت بذلك العودية البدائية لـ$\OLId{latin-ordinary}{x}\mid \OLId{latin-ordinary}{y}$.

العدد الطبيعي~$\OLId{latin-ordinary}{x}$ \emph{أولي} متى كان غير $\OLMathNumeral{0}$ وغير $\OLMathNumeral{1}$، ولم
يقبل القسمة إلا على $\OLMathNumeral{1}$ ونفسه. فشرط القواسم في العدد
الأولي أن $\OLId{latin-ordinary}{y}\mid \OLId{latin-ordinary}{x}$ يستلزم إما $\OLId{latin-ordinary}{y}=\OLMathNumeral{1}$ وإما~$\OLId{latin-ordinary}{y}=\OLId{latin-ordinary}{x}$. فإذا استبعدنا الصفر والواحد وأردنا
اختبار أولية $\OLId{latin-ordinary}{x}$، كفانا فحص $\OLId{latin-ordinary}{y}\mid \OLId{latin-ordinary}{x}$ لكل $\OLId{latin-ordinary}{y}\le \OLId{latin-ordinary}{x}$، لأن
$\OLId{latin-ordinary}{y}\nmid \OLId{latin-ordinary}{x}$ تلقائيًا إذا كان $\OLId{latin-ordinary}{y}>\OLId{latin-ordinary}{x}$. وعليه يكون تعريف العلاقة
$\fn{Prime}(\OLId{latin-ordinary}{x})$، التي تصدق إذا وفقط إذا كان $\OLId{latin-ordinary}{x}$ أوليًا، بـ
\[
\fn{Prime}(\OLId{latin-ordinary}{x}) \defiff \OLId{latin-ordinary}{x} \geq \OLMathNumeral{2} \land \bforall{\OLId{latin-ordinary}{y} \leq \OLId{latin-ordinary}{x}}{(\OLId{latin-ordinary}{y} \mid \OLId{latin-ordinary}{x} \lif \OLId{latin-ordinary}{y}
  = \OLMathNumeral{1} \lor \OLId{latin-ordinary}{y} = \OLId{latin-ordinary}{x})}
\]
فتكون عودية بدائية.

الأعداد الأولية هي $\OLMathNumeral{2}$ و$\OLMathNumeral{3}$ و$\OLMathNumeral{5}$ و$\OLMathNumeral{7}$ و$\OLMathNumeral{11}$، وهكذا. ولنسمّ الدالة
$\OLId{latin-ordinary}{p}(\OLId{latin-ordinary}{x})$ التي تعطي العدد الأولي ذا الرتبة $\OLId{latin-ordinary}{x}$ في هذه المتتالية؛ أي
$\OLId{latin-ordinary}{p}(\OLMathNumeral{0})=\OLMathNumeral{2}$ و$\OLId{latin-ordinary}{p}(\OLMathNumeral{1})=\OLMathNumeral{3}$ و$\OLId{latin-ordinary}{p}(\OLMathNumeral{2})=\OLMathNumeral{5}$، وهكذا. (وتيسيرًا للكتابة سنكتب غالبًا
$\OLId{latin-ordinary}{p}_\OLId{latin-ordinary}{x}$ بدلًا من $\OLId{latin-ordinary}{p}(\OLId{latin-ordinary}{x})$؛ فـ$\OLId{latin-ordinary}{p}_\OLMathNumeral{0}=\OLMathNumeral{2}$ و$\OLId{latin-ordinary}{p}_\OLMathNumeral{1}=\OLMathNumeral{3}$، وهكذا.)

فإن حصلت لنا دالة $\fn{nextPrime(x)}$ تعطي أول عدد أولي يجاوز~$\OLId{latin-ordinary}{x}$،
سهل تعريف $\OLId{latin-ordinary}{p}$ بالعودية البدائية:
\begin{align*}
  \OLId{latin-ordinary}{p}(\OLMathNumeral{0}) & = \OLMathNumeral{2}\\
  \OLId{latin-ordinary}{p}(\OLId{latin-ordinary}{x}+\OLMathNumeral{1}) & = \fn{nextPrime}(\OLId{latin-ordinary}{p}(\OLId{latin-ordinary}{x}))
\end{align*}
وقيمة $\fn{nextPrime}(\OLId{latin-ordinary}{x})$ هي أصغر $\OLId{latin-ordinary}{y}$ يجمع الشرطين: $\OLId{latin-ordinary}{y}>\OLId{latin-ordinary}{x}$ وكون $\OLId{latin-ordinary}{y}$ أوليًا؛ فلا يعسر
طلبها بالبحث غير المحدود. لكننا نستغني عن ذلك، ونعرفها بالتصغير
المحدود، اعتمادًا على نتيجة لإقليدس: بين $\OLId{latin-ordinary}{x}$ و
$\fact{\OLId{latin-ordinary}{x}}+\OLMathNumeral{1}$ دائمًا عدد أولي.
\[
  \fn{nextPrime(x)} =
  \bmin{\OLId{latin-ordinary}{y} \leq \fact{\OLId{latin-ordinary}{x}}+\OLMathNumeral{1}}{(\OLId{latin-ordinary}{y} > \OLId{latin-ordinary}{x} \land \fn{Prime}(\OLId{latin-ordinary}{y}))}.
\]
فتكون $\fn{nextPrime}(\OLId{latin-ordinary}{x})$، ويتبعها $\OLId{latin-ordinary}{p}(\OLId{latin-ordinary}{x})$، عوديتين بدائيتين،
وهو أقوى من مجرد قابليتهما للحساب.

(وبرهان نتيجة إقليدس موجز: إذا كان $\OLId{latin-ordinary}{x}<\OLMathNumeral{2}$، فالعدد الأولي~$\OLMathNumeral{2}$ أكبر
من~$\OLId{latin-ordinary}{x}$ ولا يجاوز $\fact{\OLId{latin-ordinary}{x}}+\OLMathNumeral{1}$. وأما إذا كان $\OLId{latin-ordinary}{x}\geq\OLMathNumeral{2}$، فليكن $\OLId{latin-ordinary}{p}_\OLId{latin-ordinary}{n}$
أكبر عدد أولي $\le \OLId{latin-ordinary}{x}$، ولنأخذ الجداء
$\OLId{latin-ordinary}{p}=\OLId{latin-ordinary}{p}_\OLMathNumeral{0}\cdot \OLId{latin-ordinary}{p}_\OLMathNumeral{1}\cdot\dots\cdot \OLId{latin-ordinary}{p}_\OLId{latin-ordinary}{n}$ لجميع الأعداد الأولية~$\le \OLId{latin-ordinary}{x}$.
إما أن يكون $\OLId{latin-ordinary}{p}+\OLMathNumeral{1}$ أوليًا، وإما أن يوجد عدد أولي بين $\OLId{latin-ordinary}{x}$ و~$\OLId{latin-ordinary}{p}+\OLMathNumeral{1}$.
وبيانه أنه إن لم يكن $\OLId{latin-ordinary}{p}+\OLMathNumeral{1}$ أوليًا، فلا بد من عدد أولي $\OLId{latin-ordinary}{q}\mid \OLId{latin-ordinary}{p}+\OLMathNumeral{1}$
بحيث $\OLId{latin-ordinary}{q}<\OLId{latin-ordinary}{p}+\OLMathNumeral{1}$. ولا يقسم أي عدد أولي $\le \OLId{latin-ordinary}{x}$ العدد $\OLId{latin-ordinary}{p}+\OLMathNumeral{1}$. (فبحسب
تعريف~$\OLId{latin-ordinary}{p}$، يقسم كل عدد أولي $\OLId{latin-ordinary}{p}_\OLId{latin-ordinary}{i}\le \OLId{latin-ordinary}{x}$ العدد~$\OLId{latin-ordinary}{p}$ بلا باق. ولذلك تكون
قسمة $\OLId{latin-ordinary}{p}+\OLMathNumeral{1}$ على كل $\OLId{latin-ordinary}{p}_\OLId{latin-ordinary}{i}\le \OLId{latin-ordinary}{x}$ ذات باق~$\OLMathNumeral{1}$، ومن ثم
$\OLId{latin-ordinary}{p}_\OLId{latin-ordinary}{i}\nmid \OLId{latin-ordinary}{p}+\OLMathNumeral{1}$.) وهكذا يكون $\OLId{latin-ordinary}{q}$ عددًا أوليًا أكبر من~$\OLId{latin-ordinary}{x}$ وأصغر من
$\OLId{latin-ordinary}{p}+\OLMathNumeral{1}$. ولأن $\OLId{latin-ordinary}{p}\le\fact{\OLId{latin-ordinary}{x}}$، يوجد عدد أولي أكبر من~$\OLId{latin-ordinary}{x}$ وأصغر من أو
مساو لـ$\fact{\OLId{latin-ordinary}{x}}+\OLMathNumeral{1}$.)

% تصحيح مصدر (C219-01): فُصلت حالتا x=0 وx=1 قبل اختيار أكبر أولي لا يجاوز x.
\textbf{تنبيه تصحيحي على الأصل.} فُصلت حالتا $\OLId{latin-ordinary}{x}=\OLMathNumeral{0}$ و$\OLId{latin-ordinary}{x}=\OLMathNumeral{1}$؛ إذ لا
يتأتى فيهما اختيار «أكبر عدد أولي لا يجاوز~$\OLId{latin-ordinary}{x}$»، ويشهد لهما العدد~$\OLMathNumeral{2}$
مباشرة.

\begin{prob}
ضع تعريفًا للقسمة الصحيحة $\OLId{latin-ordinary}{d}(\OLId{latin-ordinary}{x},\OLId{latin-ordinary}{y})$ بالتصغير المحدود.
\end{prob}

\end{document}
